\chapter*{要約}
\addcontentsline{toc}{chapter}{要約}

本論文では，線形方程式の解に対応する量子状態を生成する
HHLアルゴリズムの数学的構成を説明する．
線形方程式は，線形代数の基本問題であるとともに，
機械学習における最小二乗法にも現れる．
HHLアルゴリズムを理解するには，行列の固有値分解に加えて，
量子状態の変換と測定を組み合わせた計算を追う必要がある．
本論文の目的は，必要な量子計算の基礎を整え，
各段階の状態変化から成功時の出力を導くことである．

まず，量子状態を複素数を成分にもつベクトルで表す．
次に，量子状態の変化と測定を式で定義する．
複数の量子ビットを一つの系として扱うためにテンソル積を導入し，
基底，内積，線形写像の合成と共役転置の性質を確認する．
さらに，制御回転による振幅の変化と，測定後の状態の求め方を説明する．
次に，量子フーリエ変換，量子位相推定，振幅増幅，
行列の指数関数を取り上げ，HHLで用いる役割を明らかにする．

逆行列をもつエルミート行列$A$と，長さを1にそろえた入力状態$\ket b$に対し，
HHLでは，位相推定によって固有値を量子レジスタへ記録し，
制御回転によって固有値の逆数を振幅に取り入れる．
逆位相推定と補助量子ビットの測定を行い，
固有値の記録と各演算を誤差なく実行できると仮定し，
成功時の出力が
\[
  \ket x=\frac{A^{-1}\ket b}{\lVert A^{-1}\ket b\rVert}
\]
となることを証明する．
また，非対角な$2\times2$行列の具体例を構成し，
固有値の記録から成功確率$5/8$と，測定に成功したときの状態までを計算する．

最後に，成功確率と条件数の関係，振幅増幅の役割，
固有値の近似誤差が逆数に及ぼす影響を整理する．
計算量の評価には，入力状態の準備，制御された行列指数関数の実行，
制御回転，反復および出力の読み出しの費用が必要である．
得られる出力は解に比例する量子状態であり，
この結果だけから，任意の線形方程式の解を高速に求められるとはいえない．
このように，理想的な場合の正しさと，計算上の利点に必要な条件を区別する．

\newpage
